\appendix
\section{附录}

\subsection{附录A：特征工程流程与字段映射}

原始数据由pickdata1（资源变更记录，33字段）和pickdata2（用户行为事件记录，52字段）组成。表\ref{tab:field_map}列出了关键字段与实际数据来源的映射关系。

\begin{table}[H]
\centering
\caption{关键字段映射}
\label{tab:field_map}
\begin{tabular}{c|c|c}
\toprule
\textbf{题目需求} & \textbf{实际来源} & \textbf{提取方式} \\
\midrule
玩家基础信息 & pickdata2: register\_time, total\_pay, current\_level & 按用户聚合取最新值 \\
登录活跃 & pickdata2: dt, account\_id & 按用户$+$日期聚合 \\
付费信息 & pickdata2: total\_pay, first\_pay\_time & 取用户级max值 \\
资源消耗 & pickdata1: resource\_name, change\_type, change\_num & 筛选reduce，按资源名$+$用户聚合 \\
\bottomrule
\end{tabular}
\end{table}

特征工程流程如下：
\begin{enumerate}
    \item 分块读取pickdata1（1,931,166行）和pickdata2（4,158,533行）
    \item 按account\_id分组聚合，计算各用户的39个特征
    \item 输出player\_features.csv（2000行 $\times$ 39列），供后续各问题使用
\end{enumerate}

\subsection{附录B：策略方案完整表}

优化方案的完整8档礼包策略方案见表\ref{tab:full_strategy}。

\begin{table}[H]
\centering
\caption{优化方案策略完整表（8档礼包）}
\label{tab:full_strategy}
\resizebox{\textwidth}{!}{%
\begin{tabular}{c|c|c|c|c}
\toprule
\textbf{礼包名称} & \textbf{定价(元)} & \textbf{推送时机} & \textbf{触发条件} & \textbf{目标人群} \\
\midrule
免费挽留包 & 0 & Day 1--30 & Lost Cause象限 & 零氪浅层党、零氪流失党 \\
首充礼包 & 6 & Day 1 & Persuadable, 注册首日 & 零氪休闲党、零氪浅层党、中氪战力党 \\
新手补给 & 12 & Day 3 & Persuadable, 活跃$\geq$2天未付费 & 零氪休闲党、零氪浅层党 \\
资源补给包 & 30 & Day 7 & Persuadable, 活跃$\geq$5天或钻石$<$566 & 中氪战力党 \\
成长加速包 & 68 & Day 10 & Persuadable, 连续2天资源缺口$>$30\% & 中氪战力党 \\
战力突破包 & 128 & Day 14 & Persuadable, 等级停滞$\geq$3天 & 中氪战力党、高氪核心党 \\
至尊大礼包 & 328 & Day 21 & Persuadable高氪, 等级$\geq$10且已购$\geq$68元包 & 高氪核心党 \\
传说大礼包 & 648 & Day 28 & Persuadable高氪, 已购328元包 & 高氪核心党 \\
\bottomrule
\end{tabular}%
}
\end{table}

注：零氪用户四象限分类（零氪休闲党 Persuadable 70\%、Lost Cause 25\%；零氪浅层党 Persuadable 50\%、Lost Cause 45\%；零氪流失党 Persuadable 30\%、Lost Cause 65\%；Sleeping Dog 统一5\%），付费用户全量推送无象限区分。状态触发：生命周期$<$3天跳过$>$30元，高氪跳过$<$30元，已购$\geq$128元跳过$<$68元。

\subsection{附录C：蒙特卡洛模拟参数设置}

\begin{table}[H]
\centering
\caption{蒙特卡洛模拟关键参数}
\label{tab:mc_params}
\begin{tabular}{l|c}
\toprule
\textbf{参数} & \textbf{设定值} \\
\midrule
模拟玩家总数 $N$ & 10,000 \\
模拟重复次数 & 5,000 \\
玩家聚类数 $K$ & 6 \\
礼包数量 & 8（含免费挽留包） \\
三方案 & 基线 / 基准方案 / 优化方案 \\
留存率目标约束 & $\geq$ 10\% \\
营收优化目标 & $\max$ 总营收（目标70,000元） \\
随机种子 & 42 \\
\bottomrule
\end{tabular}
\end{table}


\subsection{附录D：资源缺口程度的定义说明}

题目要求分析"资源缺口程度"对付费金额的影响。理想的资源缺口应定义为玩家当前升级/建造所需资源与持有资源之间的差额比率：

\begin{equation}
\text{资源缺口率} = \frac{\text{升级所需资源量} - \text{当前持有资源量}}{\text{升级所需资源量}}
\end{equation}

该指标直接反映玩家为推进进度而付费的紧迫性——缺口率越高，付费动机越强。然而，现有数据仅包含资源变更记录（获取/消耗量），缺少两个关键字段：（1）各玩家当前升级/建造目标所需的资源量；（2）资源消耗的目的分类（升级消耗 vs 日常维护消耗）。因此无法精确计算上述缺口率。

本文在Hurdle模型阶段1（SHAP分析）和阶段2（Gamma GLM）中以\texttt{resource\_intensity}（日均资源消耗量 = 总消耗量 / 生命周期天数）作为资源缺口程度的代理变量。该指标衡量的是资源\textbf{消费强度}而非严格意义上的缺口，存在两类偏差：（1）高活跃度玩家的消耗量天然偏高，即使缺口率不高；（2）囤积型玩家可能持有大量资源但消耗极少，被错误标记为低缺口。建议在问题4的数据采集中补充"当前升级目标所需资源量"和"资源消耗目的标签"，以实现精确的缺口率计算。

\subsection{附录E：AI工具使用说明}

本文在建模与论文撰写过程中使用了AI辅助工具（Claude Code），主要使用方式如下：

\textbf{1. Cox预测模型是什么？}

在问题1建模初期，向AI询问Cox比例风险模型的基本原理和适用场景。AI回复要点：（1）Cox模型是半参数生存分析模型，不需指定基准风险函数的具体形式；（2）假设各协变量的风险比不随时间变化（PH假设）；（3）输出风险比$\exp(\beta)$，$\exp(\beta)<1$为保护因子，$>1$为风险因子。上述信息用于理解模型选择和论文中的方法描述。对模型的具体实现、特征选取、代码编写由本文独立完成。

\textbf{2. 游戏付费策略相关论文搜索}

在问题2和问题3的方法选择阶段，向AI询问游戏付费策略分析的参考文献。AI推荐了Ascarza等（2025）关于个性化游戏设计与用户留存的研究、Wei等（2024）关于多干预Uplift建模的论文、Belotti等（2015）关于两阶段模型（Two-Part Models）的方法论文等。参考这些文献，本文独立设计了"两阶段Hurdle模型"和"四象限Uplift框架"的建模方案，所有代码和数学推导均为独立完成。

\textbf{3. 数据解析的重要信息}

在特征工程阶段，向AI询问pickdata1和pickdata2的数据结构解析建议。AI提示注意：（1）pickdata2的register\_time字段用于确定Day1时间基准；（2）pickdata1的change\_type字段区分资源获取与消耗；（3）total\_pay以美元标价需转换为人民币；（4）连续无行为记录可判定流失。基于上述提示，本文独立完成了分块读取、特征聚合和标签构建的完整代码。

\textbf{4. 模型推广价值分析}

在论文撰写阶段，向AI询问本方法体系的潜在推广场景。AI建议可推广至SLG/卡牌/RPG类游戏、订阅制产品的用户流失分析、无A/B测试条件下的策略优化等。上述建议整合进论文的"模型评价与推广"章节，具体论证和表述由本文根据自身研究结果独立完成。

\textbf{使用方式总结：}AI工具主要用于文献检索、方法原理解释和写作建议。所有模型公式推导、代码实现、数值计算、图表生成和论文核心论证均由作者独立完成。AI生成的内容均经过作者核查和改写，确保准确性和原创性。
